\section{Three-Phase Transmission-Line Model of CHSH with Explicit Detector Physics}

\subsection{Overview}

We construct a fully causal, energy-conserving three-phase transmission-line (TL) model with angle-selective dissipative boundary conditions and evaluate CHSH correlations under two regimes:

\begin{enumerate}
    \item \textbf{Loophole-free run:} fixed integration windows, no post-selection, equal efficiencies, all trials counted.
    \item \textbf{Realistic detection run:} finite thresholds, coincidence windows, and angle-dependent detection efficiencies.
\end{enumerate}

The goal is not to evade Bell's theorem but to explicitly instantiate a classical field model in which all assumptions are transparent and auditable.

\subsection{Propagation Model}

Let $\mathbf v(x,t), \mathbf i(x,t) \in \mathbb{R}^3$ denote phase voltages and currents on a lossless three-phase transmission line of length $\ell$.

The coupled telegrapher equations are
\begin{align}
\partial_x \mathbf v(x,t) &= -\mathbf L\, \partial_t \mathbf i(x,t), \\
\partial_x \mathbf i(x,t) &= -\mathbf C\, \partial_t \mathbf v(x,t),
\end{align}
where $\mathbf L, \mathbf C$ are symmetric positive-definite per-unit-length matrices.

In the baseline model we take
\begin{equation}
\mathbf L = L' \mathbf I, \qquad \mathbf C = C' \mathbf I,
\end{equation}
yielding propagation speed
\begin{equation}
c = \frac{1}{\sqrt{L'C'}}, \qquad Z_0 = \sqrt{\frac{L'}{C'}}.
\end{equation}

Energy density and power flux are
\begin{align}
u(x,t) &= \frac{1}{2}\mathbf v^\top \mathbf C\, \mathbf v + \frac{1}{2}\mathbf i^\top \mathbf L\, \mathbf i, \\
S(x,t) &= \mathbf v^\top \mathbf i,
\end{align}
with conservation law
\begin{equation}
\partial_t u + \partial_x S = 0.
\end{equation}

Finite-speed causality follows directly from the hyperbolic PDE structure.

\subsection{Internal Polarization Space: Clarke Transform}

Define the power-invariant Clarke transform
\begin{equation}
\begin{bmatrix}
v_\alpha \\
v_\beta \\
v_0
\end{bmatrix}
= \mathbf T_C
\begin{bmatrix}
v_a \\
v_b \\
v_c
\end{bmatrix},
\end{equation}
and similarly for currents.

Balanced operation implies $v_0 \approx 0$.

Define the complex space vector
\begin{equation}
z_v = v_\alpha + i v_\beta.
\end{equation}

A detector angle $\theta$ corresponds to rotation in the $\alpha\beta$ plane:
\begin{equation}
\begin{bmatrix}
v_\alpha' \\
v_\beta'
\end{bmatrix}
=
\mathbf R(\theta)
\begin{bmatrix}
v_\alpha \\
v_\beta
\end{bmatrix},
\quad
\mathbf R(\theta)
=
\begin{bmatrix}
\cos\theta & \sin\theta \\
-\sin\theta & \cos\theta
\end{bmatrix}.
\end{equation}

Equivalently, $z_v' = e^{-i\theta} z_v$.

\subsection{Detector as Dissipative Boundary Condition}

Each port ($x=0$ or $x=\ell$) is terminated by a shunt conductance matrix $\mathbf G_{abc}(\theta)$ implementing angle-selective dissipation.

In the rotated $\alpha'\beta'0$ frame,
\begin{equation}
\mathbf G_{\alpha\beta0} =
\begin{bmatrix}
g_\parallel & 0 & 0 \\
0 & g_\perp & 0 \\
0 & 0 & g_0
\end{bmatrix}.
\end{equation}

Transforming back to phase coordinates:
\begin{equation}
\mathbf G_{abc}(\theta)
=
\mathbf T_C^\top
\begin{bmatrix}
\mathbf R(\theta)^\top
\begin{bmatrix}
g_\parallel & 0 \\
0 & g_\perp
\end{bmatrix}
\mathbf R(\theta) & \mathbf 0 \\
\mathbf 0 & g_0
\end{bmatrix}
\mathbf T_C.
\end{equation}

The boundary condition is
\begin{equation}
\mathbf i(x_{\text{port}},t) = \mathbf G_{abc}(\theta)\, \mathbf v(x_{\text{port}},t).
\end{equation}

Instantaneous absorbed power:
\begin{equation}
P_{\text{abs}}(t)
=
\mathbf v^\top \mathbf G_{abc}(\theta)\, \mathbf v
\ge 0.
\end{equation}

Channel powers:
\begin{align}
P_\parallel(t) &= g_\parallel \big(v_\alpha'(t)\big)^2, \\
P_\perp(t) &= g_\perp \big(v_\beta'(t)\big)^2.
\end{align}

\subsection{Source Model}

A wavepacket is injected at $x=\ell/2$ with internal polarization
\begin{equation}
z_s(t) = A e^{-(t/\sigma)^2} e^{i(\omega_0 t + \phi)},
\end{equation}
where $\phi$ is a per-trial random phase (hidden variable $\lambda$).

Converted to real $\alpha\beta$ components:
\begin{align}
v_{\alpha,s}(t) &= \Re[z_s(t)], \\
v_{\beta,s}(t) &= \Im[z_s(t)].
\end{align}

Phase voltages follow by inverse Clarke transform.

\subsection{Binary Outcomes}

Define fixed integration window $T$ beginning at $t_0$:
\begin{align}
E_\parallel &= \int_{t_0}^{t_0+T} P_\parallel(t)\, dt, \\
E_\perp &= \int_{t_0}^{t_0+T} P_\perp(t)\, dt.
\end{align}

Binary outcome:
\begin{equation}
A(\theta) = \operatorname{sign}(E_\parallel - E_\perp) \in \{+1,-1\}.
\end{equation}

All trials are counted (loophole-free condition).

\subsection{CHSH Evaluation}

Settings:
\begin{equation}
\theta_A \in \{0,\pi/4\}, \qquad
\theta_B \in \{\pi/8, 3\pi/8\}.
\end{equation}

Correlation:
\begin{equation}
\mathcal E(\theta_A,\theta_B)
=
\frac{1}{N}\sum_{k=1}^N
A^{(k)}(\theta_A) B^{(k)}(\theta_B).
\end{equation}

CHSH parameter:
\begin{equation}
S =
\left|
\mathcal E(\theta_A,\theta_B)
+
\mathcal E(\theta_A,\theta_B')
+
\mathcal E(\theta_A',\theta_B)
-
\mathcal E(\theta_A',\theta_B')
\right|.
\end{equation}

\subsection{Expected Outcomes}

\paragraph{Run 1 (Loophole-Free).}
With
\begin{itemize}
    \item $g_\parallel = g_\perp$,
    \item fixed $T$ and $t_0$,
    \item measurement independence,
    \item all trials counted,
\end{itemize}
the model satisfies
\begin{equation}
S \le 2,
\end{equation}
as required by Bell's theorem under local deterministic dynamics.

\paragraph{Run 2 (Realistic Detection).}
Introducing thresholds $\eta$, coincidence windows $\Delta$, or angle-dependent efficiencies may produce
\begin{equation}
S > 2
\end{equation}
via detection bias. Mapping $S(\eta,T,\Delta)$ characterizes the detection-loophole structure explicitly.

\subsection{Audit Conditions}

To ensure physical validity:
\begin{enumerate}
    \item Verify finite-speed causality (no change at remote port before $\ell/c$).
    \item Verify energy conservation:
    \[
    W_{\text{in}}(t)
    =
    \int_0^t (P_{\text{abs,A}} + P_{\text{abs,B}})\, dt
    + \int_0^\ell u(x,t)\, dx.
    \]
    \item Confirm no post-selection in Run 1.
\end{enumerate}

\subsection{Interpretation}

The model provides a concrete classical dynamical instantiation of:
\begin{itemize}
    \item global mode structure,
    \item contextual measurement via boundary conditions,
    \item local dissipative energy extraction.
\end{itemize}

It does not violate Bell's theorem under loophole-free assumptions but offers a controlled framework for analyzing how realistic detection processes can generate apparent quantum-like correlations.
